% \documentclass[a4paper,10pt]{book}

\documentclass[12pt,twoside,a4paper,english]{book}

%%%%%%%%%%%%%%%%%%%%%%%%%%%%%
%% BASIC CONFIG PARAMETERS %%
%%%%%%%%%%%%%%%%%%%%%%%%%%%%%

\newcommand{\eachPageMsg}{EMCL Thesis: WTF.}
%%%%%%%%%%%%%%%%%%%%%%%%%%%%%%%%%%%%%%%%%%%%%%%%%%%%%%%%%%%%%%%%%%%%%%%%% - July 26, 2009} % Example.

%% Macros, Packages and so on.
%%%%%%%%%%%%%%%%%%%%%%%%%%%%%%%%%%%%%%%%%%%%%%%%%%%%%%%%%%%%%%%%%%%%%%
%% Paquetes relaciones con el lenguaje y codificación general.

%% Codificación en UTF-8
% \usepackage{ucs}
\usepackage[utf8x]{inputenc}
% \usepackage[utf8]{inputenc}
% \selectlanguage{spanish} 

%% En Tetex 3, hay que usar otro formato como:
% \usepackage{ucs}
% \usepackage[utf8x]{inputenc}

%% Babel va sin opciones porque las especificamos globalmente en el
%% documento. Esto tiene la ventaja de que los demás paquetes soportan
%% español automáticamente.
\usepackage{babel}

%% Para propósitos de documentación, éste es el código antiguo para
%% seleccionar el idioma. Ya no es necesario.

%% selección del idioma
%% (ver summary.tex para activar otro idioma, ej. inglés)
%%\selectlanguage{spanish}
%%\selectlanguage{spanish} 
%% AHN: que nadie pregunte por qué
%% egallego: No, no está curado, parece ser que el problema es que hay
%% que ponerlo en sitios específicos, se ha arreglado especificando
%% las opciones correctas al paquete babel.

%%%%%%%%%%%%%%%%%%%%%%%%%%%%%%%%%%%%%%%%%%%%%%%%%%%%%%%%%%%%%%%%%%%%%%
%% Paquetes de utilidad.
%% (añadir otros paquetes necesarios para el tfc)

%% Gráficos.
%% \usepackage[final]{graphicx}
\usepackage[pdftex]{graphicx}

%% OJO!, debemos seleccionar la fuente antes del encoding, si no al
%% usar fontenc cargará las fuente EC.

%% FIXME: Aunque debería estar ya, debo elegir todavía la fuente con
%% la cual voy a imprimir el proyecto.
% \usepackage{times} % Esta no me gusta
\usepackage{palatino} % Cuidado con los < y > q no van.
% \usepackage{fourier} % Con esta no van los parenthesis grandes

%% En teoría este paquete nos proporciona una hypenation mejor.
%\usepackage[T1]{fontenc}

%% Mathematical packages.
\usepackage{amsfonts}
\usepackage{amsmath}
\usepackage{amssymb}
\usepackage{latexsym}
\usepackage{amsthm}
\usepackage{mathpazo}
\usepackage{stmaryrd}
\usepackage{mathbbol}
\usepackage[mathcal]{euscript}
% \usepackage{eufrak}

% If activated \usepackage{latexsym} 
% ERROR: LaTeX Error: Too many math alphabets used in version normal.

%% Símbolos de gramáticas. NO ACTIVAR
% \usepackage{txfonts} 
% \usepackage{yfonts}

%% Preventing figures from appearing on a page by themselves.
%% http://dcwww.camd.dtu.dk/~schiotz/comp/LatexTips/LatexTips.html
\renewcommand{\topfraction}{0.85}
\renewcommand{\textfraction}{0.1}
\renewcommand{\floatpagefraction}{0.75}


%% Mítico.
\usepackage{url}

%% Añade fixme support.
%% \usepackage{fixme}

%% Índices de contenido.
\usepackage{makeidx}
\makeindex

%% Para la bibliografía
\usepackage{natbib}
\renewcommand{\cite}[1]{\citep{#1}}

%% Arbolitos lindos
%\usepackage{pstricks,pst-node,pst-tree}
%\psset{nodesep=2pt,levelsep=40pt}

%% Formato de las cabeceras de los capítulos.
%% No instalado: 
%% \usepackage[Lenny]{fncychap}
%% Casi seguro que hay que hacerlo con titlesec.

%% Formato de la tabla de contenidos.
%% Se necesita mucho trabajo aquí todavía, entre otras cosas poner un
%% \let\chaptername\appendixname, es decir, redefinir \appendix, por
%% lo que yo creo que esto es para el final. Premature optimization is
%% the root of all evil:)
\makeatletter
\usepackage{titletoc}
\titlecontents{chapter} [6pc]
   {\addvspace{1pc}\bfseries
     \titlerule[2pt]\filright}
   {\contentslabel
      [\textsc{\chaptertitlename}\
       \thecontentslabel]{6pc}%
       \scshape}
   {\contentsmargin{0pt}}
   {\hfill\contentspage}
   [\addvspace{2pt}]
\makeatother

\usepackage{titlesec}
%% FIXME part
% \titleformat{\part}[display]
%   {\center\bfseries\Huge}
%   {\partname~\thepart}
%   {4ex}
%   {}
%   [\vspace{2ex}\titlerule]

\titleformat{\chapter}[display]
{\bfseries\scshape\LARGE}
  {\filleft\Huge\chaptertitlename~\thechapter}
  {4ex}
  {}
  [\vspace{2ex}\titlerule]

%% Hay que modificar apéndice, creo que no es demansiado difícil
% \newcommand\appendix{\par
%   \setcounter{chapter}{0}%
%   \setcounter{section}{0}%
%   \gdef\@chapapp{\appendixname}%
%   \gdef\thechapter{\@Alph\c@chapter}}

% \usepackage{appendix}
% \usepackage{soul}

%% Para q los frames se puedan dividir en páginas.
\usepackage{framed}

%% Listings.
%% Esto habrá que retocarlo un montón.

%% La opción final se sobrepone a la global de draft.
\usepackage{verbatim}
\usepackage{multicol} % Para listings con 2 cols.
% \usepackage[final]{listings}
\usepackage{listings}

%% Cuidado con la linea de literate, que contiene un hack para poder
%% poner caracteres acentuados en utf-8. Vaya chapuza de las gordas :)
% \lstset{
%   basewidth=0.5em,% OJO!! 
%   basicstyle=\fontfamily{txtt}\fontseries{m}\selectfont\small,%
%   commentstyle=\fontshape{it}\selectfont,%
%   mathescape=true,%
%   literate={->}{{$\to$}}3 {=>}{{$\Rightarrow$}}3% 
%   {-->}{{$\qquad \leadsto$}}2%
%   {á}{{\'a}}1 {é}{{\'e}}1 {í}{{\'i}}1 {ó}{{\'o}}1 {ú}{{\'u}}1, %
%   lineskip=-0.3ex,%
%   breaklines=true,%
%   frame=single,%
%   framexleftmargin=3mm,% FIXME
%   language=haskell,%
%   breakatwhitespace=true,%
%   extendedchars=true,%
%   prebreak={\mbox{\hfill$\hookleftarrow$}},%
%   morecomment=[s][\underbar]{!+}{-!} % This is for eval repr.
% }


\usepackage{charter}
\usepackage{comment}
% \usepackage{xmpincl}
\usepackage{moreverb}
\usepackage{url}
\usepackage{tabularx}
\usepackage{graphicx}
\usepackage{appendix}
\usepackage{enumerate}

% START From http://ringlord.com/publications/latex-pdf-howto/
\usepackage{color}
% \definecolor{rltred}{rgb}{0.75,0,0}
% \definecolor{rltgreen}{rgb}{0,0.5,0}
% \definecolor{rltblue}{rgb}{0,0,0.75}
% \definecolor{gris50}{rgb}{0.5,0.5,0.5}
% \definecolor{YellowOrange}{cmyk}{0,0.42,1,0}
% \definecolor{Black}{rgb}{0,0,0}
% \definecolor{White}{rgb}{1,1,1}

%\usepackage{thumbpdf}
\usepackage[pdftex,
        colorlinks=true,
        urlcolor=black,       % \href{...}{...} external (URL)
        filecolor=black,     % \href{...} local file
        linkcolor=black,       % \ref{...} and \pageref{...}
        citecolor=black,
        pdftitle={EMCL Thesis: Negative non-ground queries in Well Founded Semantics},
        pdfauthor={Victor Pablos Ceruelo},
        pdfsubject={EMCL Thesis - \today},
        pdfproducer={pdfLaTeX - \today},
        % pdfkeywords={test, testing, testable stuff},
        % pdfadjustspacing=1,
        % pagebackref,
        % pagebackref=true,
        pagebackref=false,
        pdfpagemode=UseNone,
        bookmarksopen=true]{hyperref}

% END From http://ringlord.com/publications/latex-pdf-howto/

\usepackage{footmisc}
\usepackage{longtable}
\usepackage{picinpar}
\usepackage{float}
\usepackage{hyperref}
\usepackage{graphics}    % Para los gráficos
\usepackage{tikz}
% \usepackage{lgrind}

% \floatstyle{ruled}
% \newfloat{source-code}{thp}{lop}
% \floatname{source-code}{Source Code}

%\newcounter{SourceCodeCount}
%\setcounter{SourceCodeCount}{0}
%    \stepcounter{SourceCodeCount} {\bf\small Source Code} 
%    \arabic{SourceCodeCount}: #2 \\
%    \scriptsize 

%    \line(1,0){410}
%    \line(1,0){410} \\
%    \vspace{5pt}

\newcommand{\horizontalLine}[0]{
  \begin{center}
    \begin{tikzpicture}
      \draw (0,0) -- (15,0);
    \end{tikzpicture}
  \end{center}
}
\newcommand{\horizontalLineText}[2]{
  \begin{center}
    \begin{tikzpicture}
      \draw (0,0) -- (15,0);
      \tikzstyle{every node}=[draw] 
      \node[fill=white, text centered] at (7.5, 0) {#1};
    \end{tikzpicture}
  \end{center}
}

%%%%%%%%%%%%%%%%%%%%%%%%%%%%%%%%%%%%%%%%%%%%%%%%%%%%%%%%%%%%%%%%%%%%%%
%% Sección de macros.
%% (añadir otras macros específicas para el tfc)

\def\sourcetabsize{4}
\newenvironment{sourcestyle}{\begin{scriptsize}}{\end{scriptsize}}
\def\sourceinput#1{\par\begin{sourcestyle}\verbatimtabinput[\sourcetabsize]{#1}\end{sourcestyle}\par}

% Configs varios
\urlstyle{small}

% Marcas registradas.
\def\TReg{\textsuperscript{\textregistered}}
\def\TCop{\textsuperscript{\textcopyright}}
\def\TTra{\textsuperscript{\texttrademark}}

% \TReg für ®
% \TCop für ©
% \TTra für ™

%\newcommand*\topfigrule{\hrule\kern-0.4pt\relax}
%\newcommand*\botfigrule{\hrule\kern-0.4pt\relax}

\newcommand{\labelAndCaption}[2]{
   \caption{#2}
  \label{#1}
}

\newcommand{\fixme}[1]{
  { \textbf{FIXME}~\footnote{#1} }
%  \bigskip
%  \begin{figure}[ht!]  
%    \frame{
%      \begin{minipage}[]{0.98\textwidth}
%        \smallskip
% \centering \huge
%    \newline
%        \smallskip
%      \end{minipage}
%    }
%  \end{figure}
%  \bigskip
}

\newcommand{\BoxedFormat}[1]{
%  \hfill
  \medskip
%  \vspace{0.5pt}
  \noindent
  \framebox[15.5cm]{
%  \framebox[14.7cm]{
%    \mbox
%    \hspace{0.1pt}
    \begin{minipage}[b]{0.90\textwidth} 
      % {0.98\textwidth}
      \smallskip
      #1
      \smallskip
    \end{minipage}
    \hspace{5pt}
  }
  \medskip
}


\newcommand{\InFigureBoxedFormat}[3]{
  \label{#1}
  \caption{#2}
  \BoxedFormat{#3}
  \bigskip
}

\newcommand{\ImageFormat}[4]{
  \begin{figure}
    \caption{#2}
    \label{#1}
    \includegraphics[scale=#3]{#4} 
  \end{figure}
}

\newcommand{\FileInPrologFormat}[1]{
  \lstset{language={Prolog}}
  \lstinputlisting{#1}
%    frame=lines,
%    basicstyle=\small, 
%    keywordstyle=\color{black}\bfseries\underbar,
%    stringstyle=\ttfamily, 
%    morecomment=[il]{\*}%,
%    commentstyle=\color{white}
%    deletekeywordcomment={\%\%}
%  \lstinputlisting[label=#2, caption=#3]{#1}
}

%%%%%%%%%%%%%%%%%%%%%%%%%%%%%%%%%%%%%%%%%%%%%%%%%%%%%%%%%%%%%%%%%%%%%%

%% Esto es un folclore :)
%%\usepackage{setspace}
%%\doublespacing{}

%%%%%%%%%%%%%%%%%%%%%%%%%%%%%%%%%%%%%%%%%%%%%%%%%%%%%%%%%%%%%%%%%%%%%%
%%%%%%%%%%%%%%%%%%%%%%%%%%%%%%%%%%%%%%%%%%%%%%%%%%%%%%%%%%%%%%%%%%%%%%
%%%%%%%%%%%%%%%%%%%%%%%%%%%%%%%%%%%%%%%%%%%%%%%%%%%%%%%%%%%%%%%%%%%%%%

%% New commands from ngc-2009 paper.
%% Let's see if they work ...


%\lstdefinestyle{listing}{
%  basicstyle=\ttfamily\small, 
%  commentstyle=\fontshape{it}\selectfont,%
%  breaklines=true,%
%  breakatwhitespace=true,%
%  prebreak={\mbox{\hfill$\hookleftarrow$}},%
%  numbers=none,
%  extendedchars=true,%
%  literate={->}{{$\to$}}2 {=>}{{$\Rightarrow$}}2%
%  {||}{{$\where$}}1%
%  {:-}{{$\neck~$}}2%
%  {á}{{\'a}}1 {é}{{\'e}}1 {í}{{\'i}}1 {ó}{{\'o}}1 {ú}{{\'u}}1,%
%  resetmargins=true,
%  xrightmargin=0cm,
%  mathescape=true %
%  resetmargins=true, %
%  xleftmargin=50pt
%  lineskip=-0.3ex,%
%}

\newcommand{\N}{\mathbb{N}}
\newcommand{\neck}{\leftarrow}
\newcommand{\intneg}{\mathit{intneg}}
\newcommand{\compneg}{\mathit{compneg}}
\newcommand{\CompNeg}{\mathit{CompNeg}}
\newcommand{\suc}{\mathit{succ}}
\newcommand{\eq}{\mathit{eq}}
\newcommand{\where}{[\!]} %% OJO!!
\newcommand{\df}{\mathit{def}}
\newcommand{\cdf}{\mathit{cdef}}
\newcommand{\FS}{\mathit{FS}}
\newcommand{\PS}{\mathit{PS}}
\newcommand{\n}{\mathit{neg}}
\newcommand{\Term}{\mathit{Term}}
\newcommand{\comp}{\mathit{Comp}}
\newcommand{\entails}{\models}
\newcommand{\vect}{\overline{t}}
\newcommand{\vecX}{\overline{X}}
\newcommand{\vecY}{\overline{Y}}
\newcommand{\vecx}{\overline{x}}
\newcommand{\vecy}{\overline{y}}
\newcommand{\vecz}{\overline{z}}
\newcommand{\vecb}{\overline{b}}
\newcommand{\vecB}{\overline{B}}
\newcommand{\existclos}{\boldsymbol{\exists}}
\newcommand{\CET}{\mathit{CET}}
\newcommand{\even}{\mathit{even}}
\newcommand{\then}{\rightarrow}
\newcommand{\els}{;}
\newcommand{\onlythen}{\wedge}
\newcommand{\parels}{\parallel}
\newcommand{\CD}{\mathit{CD}}
\newcommand{\NNF}{\mathit{NNF}}
\newcommand{\negp}{\mathit{neg\_p}}
\newcommand{\negrhs}{\mathit{negate\_rhs}}
\newcommand{\paratodo}{\mathit{forall}}
\newcommand{\sk}{\mathit{sk}}
% \newcommand{\next}{\mathit{next}}
\newcommand{\eval}{\mathit{eval}}
\newcommand{\Vars}{\mathit{Vars}}
\newcommand{\skolem}{\mathit{skolem}}
\newcommand{\semb}[1]{[\![{#1}]\!]}
\newcommand{\Negate}{\mathit{Negate}}
\newcommand{\depth}{\mathit{s\_depth}}
\newcommand{\vardepth}{\mathit{vardepth}}
\newcommand{\nat}{\mathit{nat}}
\newcommand{\add}{\mathit{add}}
\newcommand{\looop}{\mathit{loop}}
\newcommand{\stuck}{\Phi^{\cal H}_P}
\newcommand{\silly}{\mathit{silly}}
\newcommand{\etal}{\emph{et al}}

% VPC NO -> mathbf
\newcommand{\true}{\textbf{t}}
\newcommand{\false}{\textbf{f}}
%\newcommand{\undef}{\textbf{u}}
\newcommand{\setTrue}{\mathit{T}}
\newcommand{\setFalse}{\mathit{F}}
\newcommand{\setUndef}{\mathit{U}}
\newcommand{\interpretationI}{\mathit{I}}
\newcommand{\interpretationSetI}{\hat{\mathit{I}}}
\newcommand{\modelM}{\mathit{M}}
\newcommand{\progP}{\Pi}
% \mathit{P}}
\newcommand{\alphabetA}{\mathcal{A}}
\newcommand{\HB}{\mathcal{H}}
\newcommand{\HU}{\mathcal{U}}

%\mathbb
\newcommand{\setI}{\mathcal{I}}
\newcommand{\setJ}{\mathcal{J}}

%%%%%%%%%%%%%%%%%%%%%%%%%%%%%%%%%%%%%%%%%%%%%%%%%%%%%%%%%%%%%%%%%%%%%%
%%%%%%%%%%%%%%%%%%%%%%%%%%%%%%%%%%%%%%%%%%%%%%%%%%%%%%%%%%%%%%%%%%%%%%
%%%%%%%%%%%%%%%%%%%%%%%%%%%%%%%%%%%%%%%%%%%%%%%%%%%%%%%%%%%%%%%%%%%%%%

% From http://for.mat.bham.ac.uk/R.W.Kaye/latex/theorem/theorem.htm
% and http://www.stat.umn.edu/~charlie/amslatex.html

% Here is where all the new theorem environments are defined.
% Their use is described in the text.
\newtheorem{proposition}{Proposition}[section]
\newtheorem{theorem}{Theorem}[section]
\newtheorem{lemma}[theorem]{Lemma}%[section]
%\newtheorem*{Zorn}{Zorn's Lemma}

%\newtheoremstyle{definition}{\topsep}{\topsep}%
%{}%         Body font
%{}%         Indent amount (empty = no indent, \parindent = para indent)
%{\bfseries}% Thm head font
%{}%        Punctuation after thm head
%{0pt}%     Space after thm head (\newline = linebreak)
%{\thmname{#1}\thmnumber{ #2}\hspace{10pt}\thmnote{ #3}}%         Thm head spec

%\theoremstyle{definition}
\newtheorem{definition}{Definition}[section]

%\newtheoremstyle{corollary}{\topsep}{\topsep}%
%{}%         Body font
%{}%         Indent amount (empty = no indent, \parindent = para indent)
%{\bfseries}% Thm head font
%{}%        Punctuation after thm head
%{0pt}%     Space after thm head (\newline = linebreak)
%{\thmname{#1}\thmnumber{ #2}\hspace{10pt}\thmnote{ #3}}%         Thm head spec

%\theoremstyle{corollary}
\newtheorem{corollary}{Corollary}[section]

\newtheoremstyle{example}{\topsep}{\topsep}%
{}%         Body font
{}%         Indent amount (empty = no indent, \parindent = para indent)
{\bfseries}% Thm head font
{}%        Punctuation after thm head
{10pt}%     Space after thm head (\newline = linebreak)
{\thmname{#1}\thmnumber{ #2}\hspace{10pt}\thmnote{ #3}}%         Thm head spec

\theoremstyle{example}
\newtheorem{example}{Example}[section]

\newtheoremstyle{example*}{\topsep}{\topsep}%
{}%         Body font
{}%         Indent amount (empty = no indent, \parindent = para indent)
{\bfseries}% Thm head font
{}%        Punctuation after thm head
{10pt}%     Space after thm head (\newline = linebreak)
{\thmname{#1}\thmnote{ #3}\hspace{10pt}}%         Thm head spec

\theoremstyle{example*}
\newtheorem*{example*}{Example}

\newtheoremstyle{program}
{\topsep} %
{\topsep} %
{}%  Body font
{}%  Indent amount (empty = no indent, \parindent = para indent)
{\bfseries}% Thm head font
{}%        Punctuation after thm head
{\newline}%     Space after thm head (\newline = linebreak)
{\thmname{#1}\thmnumber{ #2}\hspace{10pt}\thmnote{ #3}}%         Thm head spec

\theoremstyle{program}
\newtheorem{program}{Program}[section]
% Do not use multicols=2. It does not work well.

\newtheoremstyle{th-head}
{\topsep} % \topsep
{} % \topsep
{\bfseries}%  Body font
{}%  Indent amount (empty = no indent, \parindent = para indent)
{\bfseries}% Thm head font
{}%        Punctuation after thm head
{0pt}%     Space after thm head (\newline = linebreak) 
{\noindent}%         Thm head spec
%\thmname{#1}  \thmnumber{ #2}\hspace{10pt}\thmnote{ #3}

\theoremstyle{th-head}
\newtheorem*{th-head}{} % [section]
% Do not use multicols=2. It does not work well.

\newtheoremstyle{remark}{\topsep}{\topsep}%
{}%         Body font
{}%         Indent amount (empty = no indent, \parindent = para indent)
{\bfseries}% Thm head font
{}%        Punctuation after thm head
{10pt}%     Space after thm head (\newline = linebreak)
{Remark.}% {\thmname{#1}\thmnumber{ #2}\hspace{10pt}\thmnote{ #3}}%         Thm head spec

\theoremstyle{remark}
\newtheorem*{remark}{Remark}
% \newtheorem{remark}{Remark}[section]

\newtheoremstyle{algorithm}
{\topsep} %
{\topsep} %
{}%  Body font
{}%  Indent amount (empty = no indent, \parindent = para indent)
{\bfseries}% Thm head font
{}%        Punctuation after thm head
{\newline}%     Space after thm head (\newline = linebreak)
{\thmname{#1}\thmnumber{ #2}\hspace{10pt}\thmnote{ #3}}%         Thm head spec

\theoremstyle{algorithm}
\newtheorem{algorithm}{Algorithm}[section]

%%%%%%%%%%%%%%%%%%%%%%%%%%%%%%%%%%%%%%%%%%%%%%%%%%%%%%%%%%%%%%%%%%%%%%
%%%%%%%%%%%%%%%%%%%%%%%%%%%%%%%%%%%%%%%%%%%%%%%%%%%%%%%%%%%%%%%%%%%%%%
%%%%%%%%%%%%%%%%%%%%%%%%%%%%%%%%%%%%%%%%%%%%%%%%%%%%%%%%%%%%%%%%%%%%%%

% How to rotate text in latex.
% From http://dmine.blogspot.com/2008/11/latex-vertical-text.html
\usepackage{rotating}

% How to justify text in latex
\usepackage{ragged2e}
% \justifying


%% Page format (beautiful headers and footers).
%% \usepackage{fancyheadings}
\usepackage{fancyhdr}

%% márgenes, tamaño del texto, etc.
\usepackage{a4wide}
%% nota: 1 inch = 2.54 centimeters
\setlength{\headheight}{15pt}
% \setlength{\topmargin}{0.99cm} % 3.5cm = 1in + 0.99cm
\setlength{\topmargin}{-0.4cm} % ~ 2 cm
% \addtolength{\topmargin}{-\headsep}
% \addtolength{\topmargin}{-\headheight}
% \setlength{\textheight}{22.5cm}
\setlength{\textheight}{23.2cm}
% \setlength{\oddsidemargin}{1.49cm} % 4cm = 1in + 1.49cm
\setlength{\oddsidemargin}{0.7cm} % ~ 3cm
% \setlength{\evensidemargin}{-0.51cm} % 2cm = 1in - 0.51cm
\setlength{\evensidemargin}{-0.5cm} % ~ 2 cm
% \setlength{\textwidth}{15cm}
\setlength{\textwidth}{15.7cm}

% Para q deje de quejarse de los saltos de linea.
% From http://newsgroups.derkeiler.com/Archive/Comp/comp.text.tex/2005-12/msg01210.html
%\setlength\parskip{.5\baselineskip
%	plus .1\baselineskip
%	minus .4\baselineskip
%}

\usepackage{etoolbox}

\fancypagestyle{plain}{%
  \fancyhead{} % get rid of the headers on plain pages
  \renewcommand{\headrulewidth}{0pt} % and the line
  \fancyfoot{} % get rid of the footers on plain pages
  \fancyfoot[C]{page \thepage} % Only the page
  % \fancyfoot[C]{\eachPageMsg{} - page \thepage} % Only the page
  \renewcommand{\footrulewidth}{0.5pt}
} 

% First you need to clear the default layout:
%\fancypagestyle{plain}{%
%  \fancyhead{}
  %\lhead{}
  %\chead{}
  %\rhead{}
%  \lfoot{}
%  \cfoot{EMCL Thesis - page \thepage}
%  \rfoot{}

  % \renewcommand{\footrulewidth}{0.5pt}

  % Do not use in this environment...
  % \fancyfoot{}  
  % \fancyfoot[C]{EMCL Thesis - page \thepage}
%}


%\pagestyle{fancy}
%\fancyhead{}
%\renewcommand\headrulewidth{.1pt}
%\fancyhead[LO,RE]{\footnotesize\sffamily\lite\leftmark}
%\fancyhead[LE,RO]{\footnotesize\sffamily\lite\itshape\rightmark}
%\fancyfoot[c]{}
%\fancyfoot[LE,RO]{\setlength{\fboxsep}{2pt}\ovalbox{\footnotesize\sffamily\thepage}}
%\fancyfoot[LO,RE]{\footnotesize\sffamily\lite\title}
%\fancypagestyle{plain}{%
%\fancyhf{}
%\fancyfoot[R]{\setlength{\fboxsep}{2pt}{\ovalbox{\footnotesize\sffamily\thepage}}}
%\fancyfoot[L]{\footnotesize\sffamily\lite\title}
%\renewcommand{\headrulewidth}{0pt} 
%}
% \pagestyle{fancy}
% This is not the way to change the format:
% \fancypagestyle{fancy}{ ... }

% There are seven letters you need to know before you can define your own header/footer:
% E: Even page
% O: Odd page
% L: Left field
% C: Center field
% R: Right field
% H: Header
% F: Footer

% Each time we select fancy format we need 
% to establish again all the parameters.
% this auxiliary macro is for that.
% This does not affect plain style pages ...
\newcommand{\resetFancyFormat}{
% First you need to clear the default layout:
\fancyhead{}
% \fancyfoot{}

% Odd (impar)
\fancyhead[OL]{\slshape \leftmark}
\fancyhead[OR]{}
% Even (par)
\fancyhead[EL]{}
\fancyhead[ER]{\slshape \rightmark}
% \fancyfoot[C]{EMCL Thesis - page \thepage} 

% 
\renewcommand{\headrulewidth}{0.5pt}
\renewcommand{\footrulewidth}{0.5pt}
}

%\newif\ifbiblio
%\renewcommand{\backmatter}{%
%  \fancyhead[C]{\ifbiblio \bibname\else\appendixname\ \thechapter\fi}%
%  \bibliotrue
%}
%\appto\appendix{\clearpage\bibliofalse}
%\preptocmd{\bibliography}{\cleardoublepage \fancyhead[C]{\bibname}}{}{}
%\preptocmd{\appendix}{\cleardoublepage \fancyhead[C]{\appendixname\ \thechapter}}{}{}
%\renewcommand{\backmatter}{
%    \fancyhdr[C]{\ifstrequal{\leftmark}{\MakeUppercase{\bibname}}
%        {\bibname}{\appendixname\ \thechapter}}
%}


% \renewcommand{\chaptermark}[1]{\markleft{#1}{}}
% \renewcommand{\sectionmark}[1]{\markright{#1}{}}
% \renewcommand{\leftmark}[0]{\chaptertitlename}
% \renewcommand{\rightmark}[0]{\sectiontitlename}
% \nouppercase (\nouppercase{\leftmark})
% \lhead{\leftmark}
% \chead{\thepage}
% \rhead{\rightmark}
% \lfoot{\today}
% \cfoot{}
% \rfoot{\thepage}
% \renewcommand{\headrulewidth}{0.4pt}
% \renewcommand{\footrulewidth}{0.4pt}

% \preptocmd{\thebibliography}{\cleardoublepage \fancyhead[C]{\bibname}}{}{}
% \preptocmd{\appendix}{\cleardoublepage \fancyhead[C]{\appendixname\ \thechapter}}{}{}

% \renewcommand{\chaptermark}[1]{%
%  \markboth{#1}{}}

%%%%%%%%%%%%%%%%%%%%%%%%%%%%%%%%%%%%%%%%%%%%%%%%%%%%%%%%%%%%%%%%%%%%%%
%%%%%%%%%%%%%%%%%%%%%%%%%%%%%%%%%%%%%%%%%%%%%%%%%%%%%%%%%%%%%%%%%%%%%%
%%%%%%%%%%%%%%%%%%%%%%%%%%%%%%%%%%%%%%%%%%%%%%%%%%%%%%%%%%%%%%%%%%%%%%

%% White pages between chapters configuration.
%% Going to odd numbered page generating a completely blankpage.
\newcommand{\clearemptydoublepage}
{\newpage{\pagestyle{empty}\cleardoublepage}}
%{\newpage{\pagestyle{empty}\clearpage}}

%%%%%%%%%%%%%%%%%%%%%%%%%%%%%%%%%%%%%%%%%%%%%%%%%%%%%%%%%%%%%%%%%%%%%%
%%%%%%%%%%%%%%%%%%%%%%%%%%%%%%%%%%%%%%%%%%%%%%%%%%%%%%%%%%%%%%%%%%%%%%
%%%%%%%%%%%%%%%%%%%%%%%%%%%%%%%%%%%%%%%%%%%%%%%%%%%%%%%%%%%%%%%%%%%%%%

% \fussy
% http://www.tac.dk/cgi-bin/info2www?(latex)\fussy
% This declaration (which is the default) makes TeX more fussy about
% line breaking. This can avoids too much space between words, but may
% produce overfull boxes.

\sloppy
% http://www.tac.dk/cgi-bin/info2www?(latex)\sloppy
% This declaration makes TeX less fussy about line breaking. This can
% prevent overfull boxes, but may leave too much space between words.
% Lasts until a `\fussy' command is issued.

%%%%%%%%%%%%%%%%%%%%%%%%%%%%%%%%%%%%%%%%%%%%%%%%%%%%%%%%%%%%%
\newcommand{\abductiveSolTestAbds}[0]
{for every term $A$, $A~\in~\Delta$, there are
  a term $B$, $~B~\in~Abdcs~$, and
  a (possible empty) substitution $~\sigma~$ 
  such that
  $~A~=~B\sigma~$ or $~A~=~(\neg B)\sigma~$}
\newcommand{\abductiveSolTestQueryInWfs}[0]
{$Q~\in~WFM~(\progP~\cup~IR~\cup~\Delta)$}
\newcommand{\abductiveSolTestBottomNotInWfs}[0]
{$\bot~\notin~WFM~(\progP~\cup~IR~\cup~\Delta)$}
%%%%%%%%%%%%%%%%%%%%%%%%%%%%%%%%%%%%%%%%%%%%%%%%%%%%%%%%%%%%%


%\documentclass{book}
\usepackage{tikz}
\usepackage{pgf}
\usepackage{subfiles}
\usepackage{amsmath, amsthm, amssymb, amsfonts}
\usepackage{listings}

\usetikzlibrary{trees,positioning,arrows}
\usetikzlibrary{decorations.pathmorphing}
\usetikzlibrary{arrows,automata,backgrounds,fit,shapes,scopes}
\newtheorem{thm}{Theorem}[section]
\newtheorem{defin}[thm]{Definition}
\newtheorem{prop}[thm]{Proposition}
\newtheorem{cor}[thm]{Corollary}
\newtheorem{ex}[thm]{Example}



\begin{document}

%%%%%%%%%%%%%%%%%%%%%%%%%%%%%%%%%%%%%%%%%%%%%%%%%%%%%%%%%%%%%%%%%%%%%%
%% Thesis initial pages:
%%%%%%%%%%%%%%%%%%%%%%%%%%%%%%%%%%%%%%%%%%%%%%%%%%%%%%%%%%%%%%%%%%%%%%

\pagestyle{empty}

%\pagestyle{empty}

\begin{titlepage}
\centering

\includegraphics[width=1\textwidth]{images/logofi.png}
\begin{center}
\line(1,0){350}\\
European Master In Computational Logic
\end{center}

\vfill

\vfill


{\Large Master Thesis}\\
\vspace{1cm}
{\huge {\sffamily\bfseries Extending Fuzzy Logic } } \\
\vspace{0.3cm}
{\huge{ \sffamily\bfseries with characteristics } } \\
\vspace{0.4cm}
{\huge{ \sffamily\bfseries Similarity and Quantification} }


\vfill

{\Large Submitted by: JINGWEI \sc{LU}}

\vspace{1cm}
{\Large Supervisor: Susana Mu\~noz Hern\'andez }\\
{\Large Co-supervisor: Victor Pablos Ceruelo } \\
\vspace{1cm}

{\Large February 2011}
\clearemptydoublepage

\end{titlepage}

% \clearpage

%%%%%%%%%%%%%%%%%%%%%%%%%%%%%%%%%%%%%%%%%%%%%%%%%%%%%%%%%%%%%%%%%
%%%%%%%%%%%%%%%%%%%%%%%%%%%%%%%%%%%%%%%%%%%%%%%%%%%%%%%%%%%%%%%%%
%%%%%%%%%%%%%%%%%%%%%%%%%%%%%%%%%%%%%%%%%%%%%%%%%%%%%%%%%%%%%%%%%

\clearemptydoublepage
\pagenumbering{arabic}
\pagestyle{plain}

%\chapter*{Abstract}

% \begin{abstract}
\documentclass[Thesis.tex]{subfiles} 
\begin{document}

%\chapter*{\centering Abstract}
\begin{abstract}
\label{chap:Abstract}

Zadeh introduced Fuzzy Logic as an extension of classical logic capable of dealing with imprecise information. It assigns a real number in $[0,1]$ to a logical statement, instead of just $0$ (false), $1$ (true).  In this work, we restrict ourselves to %two aspects interesting by themselves in Fuzzy Logic: Similarity and Quantification.
the topic: similarity between two fuzzy concepts, which is represented as fuzzy predicates in Fuzzy Logic.

Similarity between two fuzzy concepts itself is indeed a fuzzy concept. The value of similarity is
a real number between $[0,1]$. $1$ represents that two comparing concepts are completely similar, we prefer to call them identical. $0$ represents that two comparing concepts are completely different or incomparable. 

%Quantification is a fuzzy function taking fuzzy concepts as its arguments to evaluate the truth value of statement. It also can be considered as a modifier of the degree in which elements or individuals belong to a fuzzy set.

In this work, two measurements of obtaining similarity between two fuzzy predicates are proposed and presented, Interpretation Based(IBM) and Structure Based(SBM). IBM is based on set theory and corresponding relation between fuzzy set/relation and Fuzzy Logic.  In this sense, similarity relation between fuzzy predicates is achieved by their interpretations. Whilst, in BMS, fuzzy predicates are represented as predicate trees according to fuzzy rules they are involved in. The similarity between two predicates is obtained by comparison and expansion, which are taken place alternatively layer by layer till the corresponding trees are completely built.

%Quantification is defined as both second order and first order fuzzy predicate. With quantification, negation, fuzzy concept, and fuzzy rules, a complex query in Fuzzy Logic is generated to represent query and statement imprecisely in natural language. 

%Our formalization of similarity and quantification in Fuzzy Logic increase the expressivity of Fuzzy Logic by allowing us to model much more problems in real world.\\ 

\end{abstract}
\end{document}
% \end{abstract}

%%%%%%%%%%%%%%%%%%%%%%%%%%%%%%%%%%%%%%%%%%%%%%%%%%%%%%%%%%%%%%%%%
%%%%%%%%%%%%%%%%%%%%%%%%%%%%%%%%%%%%%%%%%%%%%%%%%%%%%%%%%%%%%%%%%
%%%%%%%%%%%%%%%%%%%%%%%%%%%%%%%%%%%%%%%%%%%%%%%%%%%%%%%%%%%%%%%%%

\clearemptydoublepage

%% cabeceras, pies de página y numeración hasta el primer capítulo
%% (no tocar)
% \pagenumbering{roman}
% \pagenumbering{arabic}
% \pagestyle{plain}


%%%%%%%%%%%%%%%%%%%%%%%%%%%%%%%%%%%%%%%%%%%%%%%%%%%%%%%%%%%%%%%%%%%%%%
%% Dedicatorias y agradecimientos.
%% FIXME: Debería ser un include, para poder eliminarla.

%% página con la dedicatoria
\pagestyle{empty}
\documentclass[egilmezThesis.tex]{subfiles} 
\begin{document}

\chapter*{\centering \LARGE{Acknowledgements}}
\label{chap:Acknowledgements}


\small{This work would not have been possible without the contribution of several people.}

\small{Firstly I would like to thank my supervisors Susana Mu\~noz-Hern\'andez and Victor Pablos Ceruelo for their extreme support, understanding and patience. I really appreciate the fact that they have allowed me to work on a subject, which I sincerely enjoyed and felt passionate about it through the whole process. }

\small{I would like to extend a special thanks to Victor since without his unparalleled input, this work would never be even remotely as good as it's final state. Sometimes in life lady luck smiles at you, and you meet people who are willing to help you with everything they got, regardless of the fact that they do not owe you a penny. Victor is one such unique person and I'm so glad for having the privilege of working with him.}

\small{I'm highly grateful to a list of people starting with prof. H\"{o}lldobler, Dr.Wernhard, Sylvia Wunsch, Paloma Vivas, Julia Koppenhagen and many others for organizing such an amazing program in EMCL. I have spent two wonderful years in two beautiful countries, which I will cherish for the rest of my life.}

\small{It was such a pleasure spending two years with my \textit{compadres}, Alex and Alvaro. Because of you guys I never had a dull moment during our studies.}

\small{I want to thank my dear friends Deniz, Hakan, Serdar and the \textit{projedeyim} gang for making my casual breaks from the thesis study sessions something I eagerly look forward to.}

\small{I owe a huge debt to my high-school teacher Erol Baksi for showing me the path \textit{from zero to hero}. I consider him as a great source of inspiration in my life and can only dream of being able to follow his footsteps one day. }

\small{Last but not least, I must state that I feel immensely indebted to my beloved family, but especially to my mother and my brother Selim for always being there when I needed them, for never having stopped believing in me and for showing me that for every dark night, there's a brighter day.}

\end{document}
\clearemptydoublepage

%% agradecimientos
% \input{agradecimientos}
% \clearemptydoublepage

%%%%%%%%%%%%%%%%%%%%%%%%%%%%%%%%%%%%%%%%%%%%%%%%%%%%%%%%%%%%%%%%%%%%%%
%% Contents tables
%%%%%%%%%%%%%%%%%%%%%%%%%%%%%%%%%%%%%%%%%%%%%%%%%%%%%%%%%%%%%%%%%%%%%%

%% Main contents table
\pagestyle{plain}
\addcontentsline{toc}{section}{Contents}
\printindex
\setcounter{tocdepth}{3}
\tableofcontents
\clearemptydoublepage{}

%% Figures index
\renewcommand{\thefigure}{\thechapter.\arabic{figure}}
%\renewcommand{\thefigure}{\thesection.\arabic{figure}}
%\setcounter{figure}{3}
\listoffigures
\addcontentsline{toc}{section}{Figures}
\printindex
\clearemptydoublepage{}

%% Tables index
%\addcontentsline{toc}{section}{Tables}
%\printindex
%\listoftables
%\clearemptydoublepage

% examples index (Does not work)
% \listofexamples
% \clearemptydoublepage

%% Listings index
%\addcontentsline{toc}{section}{Listings}
%\printindex
%\numberwithin{lstlisting}{section}
%\lstlistoflistings
%\clearemptydoublepage

%%%%%%%%%%%%%%%%%%%%%%%%%%%%%%%%%%%%%%%%%%%%%%%%%%%%%%%%%%%%%%%%%%%%%%
%%%%%%%%%%%%%%%%%%%%%%%%%%%%%%%%%%%%%%%%%%%%%%%%%%%%%%%%%%%%%%%%%%%%%%
%%%%%%%%%%%%%%%%%%%%%%%%%%%%%%%%%%%%%%%%%%%%%%%%%%%%%%%%%%%%%%%%%%%%%%

%% Fancy pagestyle, arabic pagenumbering.
%% Do not touch.
% \pagestyle{empty}
% \pagenumbering{arabic}
\pagestyle{fancy}
\resetFancyFormat{}

%%%%%%%%%%%%%%%%%%%%%%%%%%%%%%%%%%%%%%%%%%%%%%%%%%%%%%%%%%%%%%%%%%%%%%
%%%%%%%%%%%%%%%%%%%%%%%%%%%%%%%%%%%%%%%%%%%%%%%%%%%%%%%%%%%%%%%%%%%%%%
%%%%%%%%%%%%%%%%%%%%%%%%%%%%%%%%%%%%%%%%%%%%%%%%%%%%%%%%%%%%%%%%%%%%%%


%%%%%%%%%%%%%%%%%%%%%%%%%%%%%%%%%%%%%%%%%%%%%%%%%%%%%%%%%%%%%%%%%%%%%%
%% Thesis core: 
%%%%%%%%%%%%%%%%%%%%%%%%%%%%%%%%%%%%%%%%%%%%%%%%%%%%%%%%%%%%%%%%%%%%%%


% To include source code in thesis.
%
%\begin{figure}[h!]
%  \formatoRuby{app:fase1:migracion-respuestas}{Fase 1: Migración 1, las respuestas.}
%  \lstinputlisting{app-source-code/fase1-migracion1.rb}
%\end{figure}
%
%\\





%\newcommand{\clearemptydoublepage}
%{\newpage{\pagestyle{empty}\cleardoublepage}}

%%\pagestyle{empty}

\begin{titlepage}
\centering

\includegraphics[width=1\textwidth]{images/logofi.png}
\begin{center}
\line(1,0){350}\\
European Master In Computational Logic
\end{center}

\vfill

\vfill


{\Large Master Thesis}\\
\vspace{1cm}
{\huge {\sffamily\bfseries Extending Fuzzy Logic } } \\
\vspace{0.3cm}
{\huge{ \sffamily\bfseries with characteristics } } \\
\vspace{0.4cm}
{\huge{ \sffamily\bfseries Similarity and Quantification} }


\vfill

{\Large Submitted by: JINGWEI \sc{LU}}

\vspace{1cm}
{\Large Supervisor: Susana Mu\~noz Hern\'andez }\\
{\Large Co-supervisor: Victor Pablos Ceruelo } \\
\vspace{1cm}

{\Large February 2011}
\clearemptydoublepage

\end{titlepage}

%\newpage \thispagestyle{plain} \mbox{}
%\subfile{Abstract}

%\newpage \thispagestyle{plain} \mbox{}


%\newpage \thispagestyle{plain} \mbox{}
%\subfile{newContent/Acknowledgements}

%\newpage \thispagestyle{plain} \mbox{}
%\clearemptydoublepage
%\tableofcontents


%\newpage
\subfile{Introduction}
%\chapter{Introduction}
%\label{chap:Introduction}

\newpage \thispagestyle{plain} \mbox{}
\subfile{Preliminaries}
 

%\emptypage
\newpage
\subfile{Justification}
%\chapter{Justification}
%\label{chap:Justification}

\newpage \thispagestyle{plain} \mbox{}
\subfile{MyApproach}


\newpage \thispagestyle{plain} \mbox{}
\subfile{Error}

\subfile{Example}


\newpage \thispagestyle{plain} \mbox{}
%\chapter{Conclusion}
%\label{chap:conclusion}
\subfile{Conclusion}

%%%%%%%%%%%%%%%%%%%%%%%%%%%%%%%%%%%%%%%%%%%%%%%%%%%%%%%%%%%%%%%%%%%%%%
%% Bibliography and references.
%%%%%%%%%%%%%%%%%%%%%%%%%%%%%%%%%%%%%%%%%%%%%%%%%%%%%%%%%%%%%%%%%%%%%%

% \pagestyle{empty}
% \pagestyle{fancy}
% \resetFancyFormat{}

%% \chapter{References}
\bibliographystyle{alpha}
\phantomsection
\addcontentsline{toc}{chapter}{References}
\printindex
%%\bibliography{thesis-bib}

%\bibliography{externals/babylon/publications,externals/intneg-papers/tags/ngc/bibliography.bib}
% \input{bibliografia}
% \clearemptydoublepage{}
% \input{apendices}


%\nocite{*}
%\bibliographystyle{alpha}
\bibliography{medium-string,Thesis}

\clearemptydoublepage{}
% \clearpage 
\phantomsection
\addcontentsline{toc}{chapter}{Appendix}
\printindex
\chaptermark{Appendix}

% Odd (impar)
\fancyhead[OL]{\slshape APPENDIX}
\fancyhead[OR]{}
% Even (par)
\fancyhead[EL]{}
\fancyhead[ER]{\slshape APPENDIX}


\appendix
%\chapter*{Appendix}
\subfile{Appendix}


%%%%%%%%%


%\end{document}



%\begin{comment}
%\label{Bibliography}
%\lhead{\emph{Bibliography}}  % Change the left side page header to "Bibliography"
%\bibliographystyle{unsrtnat}  % Use the "unsrtnat" BibTeX style for formatting the Bibliography
%\bibliography{Bibliography}  % The references (bibliography) information are stored in the file named "Bibliography.bib"
%\end{comment}

\end{document}